
\section*{Appendix}
\label{sec:appendix}

Below we report on some values $f(m,n)$ and $c(m,n)$ that we have
computed by the algorithms for narrow strips described in
Subsection~\ref{subsec:narrow}. Note that Aichholzer's results referred to in
the captions of the tables have been obtained by a code that works for
\emph{general} point sets in the plane.

\medskip

\begin{table}[h]
  \centering
  {\footnotesize
    \begin{tabular}{rp{.75\textwidth}p{.12\textwidth}@{}}
      \hline
      \multicolumn{1}{c}{$n$\hspace{10mm}} &
      \multicolumn{1}{c}{$\#$ unimodular triangulations  of $P_{2,n}$} &
      \multicolumn{1}{c}{capacity} \\
      \hline
      \input 2-frozen.tex
      \hline
    \end{tabular}
  }
  \caption{Results for $m=2$ (up to $n=15$ 
     by Aichholzer~\cite{Aichholzer-WWW}).)}
  \label{tab:m2}
\end{table}


\begin{table}[h]
  \centering
  {\footnotesize 
    \begin{tabular}{rp{.75\textwidth}p{.12\textwidth}@{}}
      \hline
      \multicolumn{1}{c}{$n$\hspace{10mm}} &
      \multicolumn{1}{c}{$\#$ unimodular triangulations of $P_{3,n}$} &
      \multicolumn{1}{c}{capacity} \\
      \hline
      \input 3-frozen.tex
      \hline
    \end{tabular}
  }
  \caption{Results for $m=3$ (up to $n=10$ 
     by Aichholzer~\cite{Aichholzer-WWW}.)}
  \label{tab:m3}
\end{table}


\begin{table}[ht]
  \centering
  {\footnotesize
    \begin{sideways}
      \begin{tabular}{rrr}
        \hline
        \multicolumn{1}{c}{$n$\hspace{10mm}} &
        \multicolumn{1}{c}{$\#$ unimodular triangulations of $P_{4,n}$} &
        \multicolumn{1}{c}{capacity} \\
        \hline
        \input 4-frozen.tex
        \hline
      \end{tabular}
    \end{sideways}
  }
  \caption{Results for $m=4$ (up to $n=8$ 
      by Aichholzer~\cite{Aichholzer-WWW}.)}
  \label{tab:m4}
\end{table}



\begin{table}[ht]
  \centering
  {\footnotesize
  \begin{tabular}{rrr}
    \hline
    \multicolumn{1}{c}{$n$\hspace{10mm}} &
    \multicolumn{1}{c}{$\#$ unimodular triangulations of $P_{5,n}$} &
    \multicolumn{1}{c}{capacity} \\
    \hline
    \input 5.tex
    \hline
  \end{tabular}
  }
  \caption{Results for $m=5$ (up to $n=6$ 
    by Aichholzer~\cite{Aichholzer-WWW}.)}
  \label{tab:m5}
\end{table}


\begin{table}[ht]
  \centering
  {\footnotesize
  \begin{tabular}{rrr}
    \hline
    \multicolumn{1}{c}{$n$\hspace{10mm}} &
    \multicolumn{1}{c}{$\#$ unimodular triangulations of $P_{6,n}$} &
    \multicolumn{1}{c}{capacity} \\
    \hline
    \input 6.tex
    \hline
  \end{tabular}
  }
  \caption{Results for $m=6$.}
  \label{tab:m6}
\end{table}


%%% Local Variables:
%%% mode: latex
%%% TeX-master: "bcc_ct"
%%% End:


